%% Citas verificadas contra el PDF. Sin marcadores [VERIFY] pendientes.

\section{INTRODUCCIÓN}

El comercio electrónico peruano opera en un entorno donde la adopción del pago digital no está
acompañada por un nivel equivalente de confianza del consumidor. Un análisis comparado de
sistemas de pago en economías emergentes reporta que la falta de confianza en las
instituciones constituye una barrera clave \cite{ref26}. Solo el 57\,\% de los adultos
utiliza pagos digitales, frente al 64\,\% a nivel global. Esta desconfianza
genera un impacto directo en el mercado legítimo, ya que la intención y continuidad de compra
en canales digitales dependen fuertemente del riesgo percibido por el usuario \cite{ref21}
\cite{ref6}. La incapacidad de verificar a un vendedor desconocido antes del pago no solo
afecta a las víctimas de fraude, sino que impone una fricción comercial sobre los vendedores
honestos.

La prevención de estafas en compras digitales depende en gran medida de la capacidad del
comprador para evaluar la fiabilidad del vendedor antes de completar una transacción. Los
marcos de protección actuales se sustentan en heurísticas de uso extendido: verificación de
la reputación del vendedor, desconfianza hacia precios inverosímiles y comprobación de la
dirección del sitio web. El supuesto subyacente es que un comprador informado dispone de los
instrumentos necesarios para protegerse.

La evidencia empírica no respalda ese supuesto. Un estudio con una muestra de 820 adultos,
replicado sobre una segunda cohorte de 629, demostró que la autoconciencia sobre privacidad
digital no predijo significativamente la victimización. La sobreconfianza la incrementó
significativamente (razón de momios de 1.437; IC 95\,\% [1.234, 1.673]),
concluyendo que el conocimiento por sí solo no constituye una barrera efectiva contra el
fraude \cite{ref27}. En la misma línea, la verificación del comprador tiende a ocurrir
\emph{posterior} a que el pago ya se ha realizado \cite{ref3}.

La eficacia de estas heurísticas se ve comprometida por la pérdida de poder discriminante
de los indicadores tradicionales. Los certificados SSL, antes considerados un indicador de
fiabilidad, son obtenidos actualmente mediante servicios automatizados \cite{ref25}. La
reputación visible del vendedor ha sido demostrada como manipulable por construcción del
mercado \cite{ref29}. Las dos heurísticas más difundidas la conexión cifrada y la buena
puntuación son precisamente las que un defraudador competente puede exhibir legítimamente.

Sin embargo, existen métodos eficaces contra el fraude que operan en un plano inaccesible
para el comprador. Los sistemas desplegados por los operadores de plataformas alcanzan
desempeños elevados mediante el uso de telemetría interna \cite{ref28}. Las contramedidas
de seguridad revisadas en la literatura se implementan y financian dentro de las organizaciones
de comercio \cite{ref30}. Ninguna de estas soluciones alcanza al comprador que evalúa una
ficha en una plataforma o interactúa con un vendedor que no las ha adoptado.

El presente trabajo propone desplazar la evaluación de riesgo desde el juicio subjetivo del
comprador hacia un sistema de estimación automática que se ejecute en su navegador, sobre la
información que ya tiene disponible y antes de que se efectúe el pago. La pregunta central que
se aborda es si un conjunto de características extraíbles del lado del cliente sin acceso
privilegiado a la plataforma y sin recolección automatizada de datos resulta suficiente
para generar una advertencia útil en el punto de decisión.

Una advertencia se considera útil si cumple tres condiciones: (1) se genera antes de que el
pago se efectúe, (2) se basa en al menos cuatro características verificables, y (3)
distingue explícitamente entre riesgo detectado e información insuficiente. Esta última
condición es crítica: un sistema que presentara como seguro aquello que solo desconoce
trasladaría al comprador un riesgo que no ha evaluado.

El alcance de este sistema se delimita a la estimación probabilística del nivel de riesgo de
un vendedor mediante una escala ordinal, excluyendo la emisión de un dictamen categórico de
fraude. Esta distinción tiene dos implicaciones operativas. Por un lado, admite declaraciones
explícitas de insuficiencia de evidencia en lugar de forzar decisiones binarias. Por otro,
el sistema se evalúa en función de la mitigación de riesgo y la reducción de advertencias
infundadas, en lugar de utilizar métricas de clasificación estándar contra una verdad de
referencia inexistente en la ficha de producto.

Las contribuciones de este trabajo son tres. En primer lugar, un conjunto de ocho
características de riesgo extraíbles desde el cliente. Estas características vienen acompañadas
de una regla explícita que distingue el riesgo detectado de la evidencia insuficiente,
garantizando que la ausencia de información nunca se presente como seguridad. Esta distinción
constituye una contribución conceptual al campo: los sistemas de detección de fraude existentes
operan sobre datos completos y no abordan el problema de la información parcial que enfrenta
el comprador. En segundo lugar, una arquitectura de tres capas cuyo componente principal es una
extensión de navegador. Se presentan los diseños de los módulos clave y se entrenaron
prototipos de los clasificadores de vendedor y reseñas. En tercer lugar, un protocolo de
validación en dos fases que separa la observabilidad de las características de la eficacia del
veredicto, identificando las amenazas a la validez que condicionan las conclusiones.

Cabe precisar que la extensión de navegador y el analizador de dominio no han sido
implementados; los clasificadores han sido entrenados sobre datos sintéticos cuyas
limitaciones se analizan en la Sección~IV. El presente trabajo es, por tanto, una propuesta
arquitectónica con componentes parcialmente validados, no un sistema completo evaluado.

El resto del artículo se organiza de la siguiente manera. La Sección~II revisa el trabajo
relacionado; la Sección~III detalla la arquitectura propuesta; la Sección~IV describe el
protocolo de validación; la Sección~V presenta las conclusiones y el trabajo futuro.
